\documentclass[11pt]{article}

\usepackage[margin=1.08in]{geometry}
\usepackage{amsmath,amssymb,amsthm}
\usepackage{booktabs}
\usepackage[hidelinks]{hyperref}

\newtheorem{theorem}{Theorem}[section]
\newtheorem{proposition}[theorem]{Proposition}
\newtheorem{lemma}[theorem]{Lemma}
\newtheorem{corollary}[theorem]{Corollary}
\theoremstyle{remark}
\newtheorem{remark}[theorem]{Remark}

\newcommand{\A}{\mathbb A}
\newcommand{\C}{\mathbb C}
\newcommand{\Z}{\mathbb Z}
\newcommand{\Jac}{\operatorname{Jac}}
\newcommand{\ord}{\operatorname{ord}}
\newcommand{\lc}{\operatorname{lc}}
\newcommand{\Supp}{\operatorname{Supp}}
\newcommand{\cU}{\mathcal U}

\title{Newton-Vertex Obstructions for Polynomial-Graph Descents\\
of the Gallagher Weighted Lift}
\author{Atharva Vaidya}
\date{26 July 2026}

\begin{document}
\maketitle

\begin{abstract}
We study polynomial-plane restrictions of the Gallagher weighted lift
\(F=(A,B,C):\A^3\to\A^3\).  Let \(\cU(A,B,C)\) be its exact
second-subduction coordinate.  We prove that, on every polynomial graph
\(z=g(x,y)\),
\[
 \Jac_{x,y}\bigl(\cU+R,Q\bigr)\notin\C^\times
\]
for every target polynomial \(R(A,B,C)\) of degree at most two and every
nonconstant, not necessarily homogeneous, \(Q(A,B,C)\) of degree at most
eleven.  More generally the conclusion holds in arbitrary target degree
provided no two support exponents of \(Q\) differ by a nonzero multiple of
\((5,-6,4)\).  Its primitive positive and negative parts give the first
possible collision, \(A^5C^4\) and \(B^6\), in degrees nine and six.
We resolve this collision and all its multiples through degree eleven by an
exact triangular basis change through three cusp subductions.

The proof uses the globally earliest Newton vertex, rather than an
isolated homogeneous face.  At a regular source boundary an exact
seed gap protects either the leading Wronskian or its first
nonmonomial deviation.  At a pole boundary, cancellation at
\(u=1\) is incompatible with the leading coefficient at
\(u=\infty\).  At degree eleven six two-dimensional polar groups remain.
Four are immediately nonresonant; the two resonant characteristic rays
are eliminated by exact local Laurent-jet determinants at \(u=1\).
The determinants depend only on the graph's forced 1-jet and are
independent of every higher zero-sector coefficient.  The finite
seed-support assertions are checked by exact rational arithmetic in the
accompanying code.

This result excludes a broad descent architecture for a
three-dimensional Keller counterexample.  It neither proves nor
disproves the two-dimensional Jacobian conjecture.
\end{abstract}

\section{The weighted lift and the main theorem}

The plane Jacobian conjecture and its standard reduction theory are
surveyed in~\cite{BassConnellWright,vanDenEssen}.  The construction
studied here is instead a descent problem arising from Gallagher's
three-dimensional weighted-lift family~\cite{Gallagher}; none of the
arguments below changes the scope of the classical reductions.  A
targeted search conducted at the date of this draft found no statement
matching Theorem~\ref{thm:main}; this is a cautious literature report,
not a claim of priority.

Put
\[
 r(w)=w(w+1)(w-3)(w-4)(w-5)(w+4),\qquad
 H(w)=-\frac{r(w)}{92},
\]
and define
\[
 p(w)=H'(w)-H'(0),\qquad q(w)=wH'(w)-H(w).
\]
Their leading coefficients are
\[
 p_5=-\frac3{46},\qquad q_6=-\frac5{92}.
\]
In source coordinates put
\[
 u=1+xy,\qquad
 \gamma=1-\frac{57}{34}xy+x^2z,\qquad
 w=u\gamma .
\]
The weighted-lift coordinates have the exact presentations
\[
 A=\frac{q(w)+w\gamma}{x^2\gamma^2},\qquad
 B=\frac{p(w)+\gamma}{x\gamma},\qquad
 C=x\gamma .
\tag{1}
\]
The engineered divisibilities make the three expressions in (1)
polynomials in \(x,y,z\).  Direct differentiation gives
\[
\det\frac{\partial(A,B,C)}{\partial(x,y,z)}=1.
\]
This identity is checked from the explicit seed by
\path{verify_weighted_lift_degree_six_graph_projection_no_go.py} and
is also part of Gallagher's source construction~\cite{Gallagher}.

We record the target polynomial used as the first coordinate.  Let
\[
\begin{aligned}
t_1&=-\frac{8505}{18948593792},&
t_2&=-\frac{79515}{151588750336},\\
t_3&=\frac{1587485}{303177500672},&
t_4&=-\frac{1079635}{303177500672},
\end{aligned}
\]
and set
\[
\begin{aligned}
T={}&p_5^6A^5C^4-q_6^5B^6
 +t_1A^4BC^3+t_2A^3B^2C^2\\
&+t_3A^2B^3C+t_4AB^4.
\end{aligned}
\tag{2}
\]
Write
\[
\lambda_0=
\frac{258473511}{31221670147323559936}
\]
and
\[
\begin{aligned}
s_0&=-\frac{10204102899}{11489574614215070056448},\\
s_1&=\frac{5559334317}{249773361178588479488},\\
s_2&=-\frac{1028649031155}{11489574614215070056448},\\
s_3&=\frac{16337519127}{359049206694220939264}.
\end{aligned}
\]
The exact second subduction is
\[
\cU=p_5^5CT-\lambda_0B^5+
s_0A^4C^3+s_1A^3BC^2+s_2A^2B^2C+s_3AB^3.
\tag{3}
\]

Adding a constant to \(Q\) does not change its Jacobian with any first
coordinate.  We therefore replace \(Q\) by \(Q-Q(0)\) throughout.  For
a polynomial
\[
Q-Q(0)=
\sum_{\substack{a,b,c\ge0\\a+b+c>0}}
\lambda_{abc}A^aB^bC^c,
\]
its support is the set of exponent triples with
\(\lambda_{abc}\ne0\).  Call the support \emph{cusp-free} if no two
of its elements differ by a nonzero multiple of \((5,-6,4)\).

\begin{theorem}[Nonhomogeneous Newton-vertex obstruction]
\label{thm:main}
Let \(g\in\C[x,y]\), and restrict (1) to the graph \(z=g(x,y)\).
If \(Q\in\C[A,B,C]\setminus\C\) has cusp-free support and
\(\deg R\le2\), then
\[
\Jac_{x,y}\bigl(\cU(A,B,C)+R(A,B,C),Q(A,B,C)\bigr)
\notin\C^\times .
\tag{4}
\]
\end{theorem}

\begin{corollary}[Automatic cusp-free range through degree eight]
\label{cor:eight}
The conclusion of Theorem~\ref{thm:main} holds for every nonconstant
\(Q\) of total target degree at most eight.  No homogeneity assumption
is imposed.
\end{corollary}

\begin{theorem}[Arbitrary targets through degree ten]
\label{thm:ten}
Let \(g\in\C[x,y]\), restrict (1) to \(z=g(x,y)\), and let
\(\deg R\le2\).  For every nonconstant \(Q\in\C[A,B,C]\) of total
target degree at most ten,
\[
\Jac_{x,y}\bigl(\cU+R,Q\bigr)\notin\C^\times .
\]
\end{theorem}

\begin{theorem}[Arbitrary targets through degree eleven]
\label{thm:eleven}
Under the hypotheses of Theorem~\ref{thm:ten}, the same conclusion holds
for every nonconstant \(Q\in\C[A,B,C]\) of total target degree at most
eleven.
\end{theorem}

\begin{corollary}[Homogeneous targets]
\label{cor:homogeneous}
The conclusion of Theorem~\ref{thm:main} holds for every nonzero
homogeneous target \(Q_n(A,B,C)\), in every degree \(n\), and every
\(\deg R\le2\).
\end{corollary}

\section{Exact boundary templates}

After restricting to \(z=g(x,y)\), substitute
\[
y=\frac{u-1}{x}.
\]
Then
\[
\gamma
=1-\frac{57}{34}(u-1)
 +x^2g\left(x,\frac{u-1}{x}\right)
=\sum_{j=\rho}^{M}x^jf_j(u),
\qquad f_\rho\ne0.
\tag{5}
\]
Necessarily \(\rho\le0\).  If there is no negative power, the
coefficient of \(x^0\) takes the value \(1\) at \(u=1\).

Let \(\mathcal N(w,\gamma)\) be the numerator obtained from (3) after
using (1) and clearing \(x^5\gamma^5\).  Define
\[
\begin{aligned}
\mathsf A(u,\gamma)
 &=\frac{q(u\gamma)+u\gamma^2}{\gamma^2},\\
\mathsf B(u,\gamma)
 &=\frac{p(u\gamma)+\gamma}{\gamma},\\
\mathsf H(u,\gamma)
 &=\frac{\mathcal N(u\gamma,\gamma)}{\gamma^5}.
\end{aligned}
\tag{6}
\]
Thus
\[
A=x^{-2}\mathsf A,\qquad
B=x^{-1}\mathsf B,\qquad
C=x\gamma,\qquad
\cU=x^{-5}\mathsf H.
\tag{7}
\]

\begin{lemma}[Exact seed support]
\label{lem:seed}
The templates in (6) are polynomials and have unique top
\(\gamma\)-degree terms
\[
\begin{aligned}
\mathsf A&=q_6u^6\gamma^4+\text{terms of lower \(\gamma\)-degree},\\
\mathsf B&=p_5u^5\gamma^4+\text{terms of lower \(\gamma\)-degree},\\
\mathsf H&=\theta u^{20}\gamma^{17}
  +\text{terms of lower \(\gamma\)-degree},
\end{aligned}
\tag{8}
\]
where
\[
\theta=
-\frac{36905625}{6077984970919772059860992}\ne0.
\]
If \(\rho=0\), put \(f=f_0\) and \(L=\deg f\).  Then \(L\ge1\),
and every non-top seed term in (8), after \(\gamma=f\), loses at
least \(L+1\) in \(u\)-degree.  If \(\rho<0\), every non-top seed
term has strictly later \(x\)-order.
\end{lemma}

\begin{proof}
The assertions follow by exact expansion of the degree-six seed and
the two target subductions (2)--(3).  For \(\mathsf A,\mathsf B\),
every non-top term loses at least one \(u\)-power and one
\(\gamma\)-power.  For \(\mathsf H\), the only additional edge
possibility loses two \(\gamma\)-powers and no \(u\)-power.  The
regular loss is therefore at least
\(\min\{L+1,2L\}=L+1\).  When \(\rho<0\), losing a
\(\gamma\)-power strictly raises the \(x\)-order.

The accompanying verifier expands all 77 nonzero terms of
\(\mathcal N\) over \(\mathbb Q\), checks the divisibilities and
support inequalities, and checks the displayed value of \(\theta\).
No floating-point calculation is used.
\end{proof}

We repeatedly use the elementary identity
\[
\Jac_{x,y}\bigl(x^\alpha H(u),x^\beta V(u)\bigr)
=x^{\alpha+\beta}
\bigl(\alpha HV'-\beta H'V\bigr),
\tag{9}
\]
which follows from \(u=1+xy\).

\section{Regular boundaries}

Assume \(\rho=0\).  Then
\[
f(1)=1,\qquad f'(1)=-\frac{57}{34},\qquad L=\deg f\ge1.
\tag{10}
\]
For an exponent triple \(m=(a,b,c)\) define
\[
\beta=-2a-b+c,\qquad h=6a+5b,\qquad q=4a+4b+c.
\tag{11}
\]
Notice that \(q=h+\beta\).  The top boundary term of
\(A^aB^bC^c\) is
\[
q_6^ap_5^b x^\beta u^hf^q.
\tag{12}
\]

Let \(\beta_0\) be the minimum \(\beta\) in \(\Supp Q\), and let
\(h_0\) be the maximum \(h\) among terms with \(\beta=\beta_0\).
The exponent triples at this vertex all have the same \(q=h_0+\beta_0\).
Their aggregate coefficient is
\[
\Lambda_0=
\sum_{\substack{\lambda_{abc}\ne0\\
\beta=\beta_0,\ h=h_0}}
\lambda_{abc}q_6^ap_5^b.
\tag{13}
\]

\begin{proposition}[Regular vertex]
\label{prop:regular}
If \(\Lambda_0\ne0\), then
\(\Jac(\cU,Q)\notin\C^\times\).
\end{proposition}

\begin{proof}
On a fixed \(\beta\)-face the \(u\)-degree of (12) is
\[
h+qL=h+(h+\beta)L.
\tag{14}
\]
Decreasing \(h\) by \(r\ge1\) therefore costs exactly
\(r(L+1)\).  Lemma~\ref{lem:seed} gives the same protected gap for
non-top seed terms.  Consequently (13) supplies a unique nonzero
aggregate at the globally earliest Jacobian order
\(x^{\beta_0-5}\).

For a vertex exponent triple, the coefficient Wronskian is
\[
\mathcal W=-5\mathsf H\mathsf V'
-\beta_0\mathsf H'\mathsf V.
\tag{15}
\]
Its leading coefficient is a nonzero seed-and-target scalar times
\[
B_*L+C_*,
\tag{16}
\]
where
\[
B_*=-5q-17\beta_0=14a-3b-22c,
\quad
C_*=-5h-20\beta_0=10a-5b-20c.
\tag{17}
\]
If (16) is nonzero, then \(\mathcal W\ne0\).

Suppose (16) vanishes.  The identity
\[
17C_*-15B_*=-10q
\tag{18}
\]
implies \(B_*\ne0\).  The jets (10) rule out \(f=du^L\), so
\[
f=du^L+eu^{L-\delta}+\cdots,
\qquad de\ne0,\quad1\le\delta\le L.
\tag{19}
\]
After the leading cancellation, the first surviving coefficient in
(15) is a nonzero multiple of
\[
-B_*\delta e.
\tag{20}
\]
It occurs at degree drop at most \(L\), before any other target
vertex or non-top seed term can enter.

If \(\beta_0-5\ne0\), the nonzero earliest \(x\)-term already
contradicts a constant Jacobian.  If \(\beta_0-5=0\), the surviving
coefficient still has positive \(u\)-degree.  Before cancellation its
degree is
\[
20+17L+h_0+qL-1>0,
\]
and (20) lowers it by only \(\delta\le L\).  Thus it cannot be a
nonzero constant.
\end{proof}

\section{Pole boundaries}

Assume \(\rho<0\) and put \(f=f_\rho\).  A graph monomial
\(x^ry^s\) can contribute to \(f\) only when \(r+2-s=\rho\);
hence
\[
f(u)=\sum_r c_r(u-1)^{r+2-\rho}.
\]
Consequently
\[
\nu:=\ord_{u=1}f\ge2-\rho,\qquad
L:=\deg f\ge2-\rho,\qquad
L+\rho+1\ge3.
\tag{21}
\]
The leading boundary terms are
\[
\begin{aligned}
\cU&=x^\alpha\theta u^{20}f^{17}+\cdots,
&\alpha&=-5+17\rho,\\
A&=x^{-2+4\rho}q_6u^6f^4+\cdots,\\
B&=x^{-1+4\rho}p_5u^5f^4+\cdots,\\
C&=x^{1+\rho}f+\cdots.
\end{aligned}
\tag{22}
\]

For a target exponent triple retain \(q\) from (11) and put
\[
E=\beta+\rho q.
\tag{23}
\]
Let \(E_0\) be the minimum \(E\) in the support.  At the
globally earliest face, group the top target terms by \(q\):
\[
\Lambda_q=
\sum_{\substack{\lambda_{abc}\ne0\\E=E_0,\ q(a,b,c)=q}}
\lambda_{abc}q_6^ap_5^b.
\tag{24}
\]
Let \(q_{\min}\) and \(q_{\max}\) be the smallest and largest
indices with \(\Lambda_q\ne0\).

\begin{proposition}[Polar endpoints]
\label{prop:polar}
If some \(\Lambda_q\ne0\), then
\(\Jac(\cU,Q)\notin\C^\times\).
\end{proposition}

\begin{proof}
At fixed \(E_0\), the \(q\)-group is a nonzero multiple of
\[
u^sf^q,\qquad s=(1+\rho)q-E_0.
\tag{25}
\]
At \(u=1\), the \(q_{\min}\)-group uniquely has least vanishing
order.  Its \(f'\)-coefficient in the Wronskian is
\[
\alpha q_{\min}-17E_0.
\tag{26}
\]
If this is nonzero, the resulting term has order
\[
(17+q_{\min})\nu-1>0
\tag{27}
\]
at \(u=1\), and every larger \(q\) has strictly greater order.

Suppose instead that (26) vanishes, so
\[
E_0=\frac{\alpha q_{\min}}{17}.
\tag{28}
\]
At infinity the \(q_{\max}\)-group uniquely dominates because
\(L+\rho+1>0\).  Its leading coefficient is
\[
\alpha\left(
(L+\rho+1)(q_{\max}-q_{\min})
+\frac{2q_{\min}}{17}
\right),
\tag{29}
\]
which is nonzero.  Its degree is
\[
20+17L+(L+\rho+1)q_{\max}-E_0-1>0.
\tag{30}
\]
Here \(E_0\le0\): for \(\rho\le-1\), the coefficients of
\(a,b,c\) in \(E\) are respectively
\(-2+4\rho,-1+4\rho,1+\rho\), all nonpositive.
Every non-top seed term has later \(x\)-order by
Lemma~\ref{lem:seed}.  Hence the earliest coefficient is nonzero
and, if its \(x\)-order is zero, is still nonconstant by (27) or
(30).
\end{proof}

\section{The cusp kernel and the degree-eight cutoff}

\begin{lemma}[Primitive collision direction]
\label{lem:kernel}
Two exponent triples have the same pair \((\beta,h)\) if and only if
their difference is
\[
k(5,-6,4),\qquad k\in\Z.
\tag{31}
\]
For every fixed \(\rho<0\), equality of \((E,q)\) has the same
kernel.
\end{lemma}

\begin{proof}
The integer kernel of
\[
\begin{pmatrix}-2&-1&1\\6&5&0\end{pmatrix}
\]
is generated primitively by \((5,-6,4)\).  Replacing the first row
by \(E=\beta+\rho q\) and the second by \(q=(4,4,1)\) leaves the
same kernel.
\end{proof}

\begin{proof}[Proof of Theorem~\ref{thm:main}]
At a regular boundary, cusp-freeness makes the vertex in (13) a
single nonzero target coefficient.  Proposition~\ref{prop:regular}
applies.  At a pole boundary, Lemma~\ref{lem:kernel} makes every
nonempty aggregate in (24) nonzero, so
Proposition~\ref{prop:polar} applies.

It remains to include \(R\).  At a regular boundary, every
nonconstant monomial of target degree at most two has \(x\)-order at
least \(-4>-5=\ord_x\cU\).  At a pole boundary such a monomial has
\(\beta\ge-4\) and \(q\le8\), whence
\[
E_R-\alpha=(\beta+5)+\rho(q-17)>0.
\tag{32}
\]
Thus \(\Jac(R,Q)\) occurs strictly after the protected earliest
coefficient of \(\Jac(\cU,Q)\).
\end{proof}

\begin{proof}[Proof of Corollary~\ref{cor:eight}]
Suppose two distinct nonnegative exponent triples differ by
\(k(5,-6,4)\), with \(k>0\) after exchanging them.  The triple on
the negative-\(B\) side has \(b\ge6k\), so its total target degree
is at least \(6k\).  The other triple has total degree larger by
\(3k\), hence at least \(9k\).  Such a pair cannot occur when all
monomials have degree at most eight.
\end{proof}

\begin{proof}[Proof of Corollary~\ref{cor:homogeneous}]
The total degree changes by \(5-6+4=3\) along the primitive kernel.
Distinct monomials in one homogeneous target therefore cannot differ
by a nonzero multiple of (31).
\end{proof}

\section{Exact closure through degree ten}

The primitive positive and negative parts of (31) give the first
collision
\[
A^5C^4\quad\text{and}\quad B^6.
\tag{33}
\]
Through degree ten, this collision and its multiples give exactly
\[
\begin{array}{c|c}
B^6&A^5C^4\\
AB^6&A^6C^4\\
B^7&A^5BC^4\\
B^6C&A^5C^5.
\end{array}
\tag{34}
\]
The polynomial \(T\) in (2) replaces the upper monomial in the first
row, and \(AT\) does the same in the second row.  To handle the final
two rows, define
\[
\begin{aligned}
W={}&q_6^5BT-p_5\lambda_0A^5C^3
-\frac{1289186932405}{367666387654882241806336}A^4BC^2\\
&+\frac{967724214885}{183833193827441120903168}A^3B^2C\\
&+\frac{280059055635}{367666387654882241806336}A^2B^3.
\end{aligned}
\tag{35}
\]
The coordinate \(\cU\) in (3) replaces \(CT\), while \(W\) replaces
\(BT\).

\begin{lemma}[Exact degree-ten basis]
\label{lem:basis-ten}
Remove
\[
(5,0,4),\quad(6,0,4),\quad(5,1,4),\quad(5,0,5)
\tag{36}
\]
from the ordinary exponent triples of total degree at most ten.
Adjoining \(T,AT,W,\cU\) gives a basis of the same polynomial space.
Relative to the four removed monomials, the change-of-basis diagonal is
\[
p_5^6,\qquad p_5^6,\qquad q_6^5p_5^6,\qquad p_5^{11},
\tag{37}
\]
and is nonzero.
\end{lemma}

\begin{proof}
Lemma~\ref{lem:kernel} and the degree bound give exactly the four pairs
in (34).  Multiplying (2) by \(1,A,B,C\) first gives a triangular
replacement by \(T,AT,BT,CT\).  Equations (3) and (35) then replace
\(CT,BT\) by \(\cU,W\), with the diagonal (37).  Exact expansion over
\(\mathbb Q\) checks that the displayed target formulas agree with the
boundary builders.
\end{proof}

We next record the boundary data of the three augmented atoms that remain
after the first-coordinate reduction.  At a regular boundary their exact
tops are
\[
\begin{array}{c|c|c}
\text{atom}&\beta&u^hf^q\\ \hline
T&-6&\lambda_0u^{25}f^{19}\\
AT&-8&q_6\lambda_0u^{31}f^{23},
\end{array}
\tag{38}
\]
up to the nonzero seed factors already displayed.  Every non-top term
in these two rows loses at least one \(u\)-power and one \(f\)-power,
hence at least \(L+1\) in regular degree, and is later at a pole.
The exact Pareto frontier of \(W\) is
\[
\kappa_{19}u^{26}f^{19}
+\kappa_{20}u^{25}f^{20},
\qquad \kappa_{19}\kappa_{20}\ne0.
\tag{39}
\]
The second term is the unique polar top.  It is also the unique regular
top for \(L\ge2\).  For \(L=1\), the fixed jets force
\[
f(u)=\frac{91}{34}-\frac{57}{34}u,
\tag{40}
\]
and the two frontier terms combine to a nonzero coefficient of degree
\(45\).

\begin{proof}[Proof of Theorem~\ref{thm:ten}]
By Lemma~\ref{lem:basis-ten}, write
\[
Q=c+\mu\cU+\lambda_TT+\lambda_{AT}AT+\lambda_WW
+\sum\nolimits'\lambda_{abc}A^aB^bC^c,
\tag{41}
\]
where the prime omits (36).  Put
\[
P=\cU+R,\qquad Q^\flat=Q-\mu P.
\tag{42}
\]
Then \(\Jac(P,Q^\flat)=\Jac(P,Q)\), the coefficient of \(\cU\) in
\(Q^\flat\) is zero, and only ordinary coefficients of degree at most
two have changed.  If \(Q^\flat\) is constant, the Jacobian is zero.
Assume henceforth that it is nonconstant.

At a regular boundary, the retained ordinary \((\beta,h)\)-labels are
pairwise distinct, and neither \((-6,25)\) nor \((-8,31)\) is ordinary.
For \(L\ge2\), equality between the \(W\)-degree \(25+20L\) and an
ordinary atom of order \(\beta=-7\) would imply
\[
(L+1)(h-27)=-2,
\tag{43}
\]
which is impossible.  For \(L=1\), equality would require \(h=26\);
the equations \(6a+5b=26\), \(c=-7+2a+b\ge0\) have no nonnegative
solution.  Thus the globally earliest regular vertex is a singleton.

The leading Wronskian characteristics of the new atoms are
\[
\begin{array}{c|c}
T&7L-5\\
AT&21L+5\\
W&19L+15\quad(L\ge2),
\end{array}
\tag{44}
\]
and the \(L=1\) characteristic of \(W\) is \(34\).  All are nonzero.
For an ordinary atom whose top characteristic vanishes, the deviation
argument in Proposition~\ref{prop:regular} produces a nonzero coefficient
within degree drop at most \(L\).  The only possible new interference
could come from \(W\) on the \(\beta=-7\) face.  If an ordinary atom of
height \(h\) lies above \(W\), its degree excess is
\[
(L+1)(h-27)+2.
\tag{45}
\]
For \(h\ge28\) this is greater than \(L\); the remaining leading case
\(h=27\) is \(A^2B^3\), whose characteristic is \(19L+5\ne0\).
Hence the earliest regular Wronskian is nonzero.

At a pole, put
\[
\alpha=-5+17\rho,\qquad E=\beta+\rho q.
\tag{46}
\]
The retained ordinary polar labels, together with \(T\) and \(AT\), are
pairwise distinct.  The sole collision after adjoining \(W\) is
\[
(\beta_W,q_W)=(-7,20)
=(\beta_{A^2B^3},q_{A^2B^3}).
\tag{47}
\]
At the globally earliest \(E\), choose the least nonzero \(q\)-group.
If it is a singleton with top \(u^hf^q\), its Wronskian, up to a
nonzero monomial factor, is
\[
(\alpha q-17E)\,u f'
+(\alpha h-20E)\,f.
\tag{48}
\]
If the \(f'\)-coefficient is nonzero, this group has vanishing order
\((17+q)\nu-1\) at \(u=1\), before every larger \(q\)-group.  If it
vanishes, then \(\beta=-5q/17\), and the \(f\)-coefficient is
\[
\alpha h-20E=2E\ne0.
\tag{49}
\]
Its order is \((17+q)\nu\), still before the next group because
\(\nu\ge3\).

It remains to treat the exceptional two-dimensional group.  Its exact
polar top is
\[
f^{20}P(u),\qquad P(u)=a u^{27}+b u^{25},
\qquad(a,b)\ne(0,0).
\tag{50}
\]
Writing \(E=-7+20\rho\), direct differentiation gives, again up to a
nonzero monomial factor,
\[
19u f'P+f\bigl(\alpha uP'-20EP\bigr).
\tag{51}
\]
If \(P(1)\ne0\), its first coefficient has order \(37\nu-1\) and is
nonzero.  If \(P(1)=0\), then \(P=a u^{25}(u^2-1)\) has a simple zero,
and the coefficient at order \(37\nu\) is proportional to
\[
(19\nu+\alpha)P'(1).
\tag{52}
\]
The graph inequality \(\nu\ge2-\rho\) gives
\[
19\nu+\alpha\ge33-2\rho>0.
\tag{53}
\]
Every group with \(q\ge21\) starts later, so cancellation is impossible.

Finally, every nonconstant monomial of \(R\) is later than \(\cU\):
its regular order is at least \(-4>-5\), while at a pole
\[
E_R-\alpha=(\beta+5)+\rho(q-17)>0
\tag{54}
\]
because \(\beta\ge-4\), \(q\le8\), and \(\rho<0\).  Thus
\(\Jac(R,Q^\flat)\) cannot enter the isolated coefficient of
\(\Jac(\cU,Q^\flat)\).  If that coefficient has \(x\)-order zero, it
still has positive \(u\)-degree at a regular boundary or positive
vanishing order at a pole.  It is therefore not a nonzero constant.
\end{proof}

\begin{remark}[No SAGBI termination claim]
The generator \(W\) is forced: its two frontier vectors do not belong to
the semigroup generated by \(A,B,C,T,\cU\).  Later exact subductions
produce further projected-new values.  The proof above uses only the
finite vector space through degree ten; it asserts neither finite
termination nor infinitude of the full SAGBI completion.
\end{remark}

\section{Exact closure through degree eleven}

Degree eleven contains ten copies of the primitive collision (33), one
for each monomial of degree at most two:
\[
B^6M\sim A^5C^4M,\qquad \deg M\le2.
\tag{55}
\]
The four degree-ten replacements extend without a new subduction
generator.

\begin{lemma}[Exact degree-eleven basis]
\label{lem:basis-eleven}
Remove \(A^5C^4M\), for all monomials \(M\) of degree at most two, from
the ordinary monomial basis of target degree at most eleven.  The ten
polynomials
\[
T,\ AT,\ W,\ \cU,\ A^2T,\ AW,\ A\cU,\ BW,\ B\cU,\ C\cU
\tag{56}
\]
replace them and give a basis of the same vector space.
\end{lemma}

\begin{proof}
The collision kernel is \(\Z(5,-6,4)\), so (55) is the complete list.
Multiplication of \(T\) by
\[
1,\ A,\ B,\ C,\ A^2,\ AB,\ AC,\ B^2,\ BC,\ C^2
\]
first gives a triangular replacement.  Equations (3) and (35) then
replace the corresponding \(CT\) and \(BT\) multiples by the entries
of (56).  The exact coefficient matrix on the ten removed monomials has
nonzero determinant.  The accompanying verifier constructs this matrix
over \(\mathbb Q\).
\end{proof}

As before, write \(P=\cU+R_{\le2}\).  Subtracting the scalar
\(\cU\)-coefficient of a target modulo \(P\) changes only ordinary
coefficients of degree at most two and preserves the Jacobian.  The
remaining augmented atoms are
\[
\mathcal A_{11}=
\{T,AT,A^2T,W,AW,BW,A\cU,B\cU,C\cU\}.
\tag{57}
\]
Their exact Pareto frontiers, recorded as \((\beta,h,q)\), are
\[
\begin{array}{c|c}
\text{atom}&\text{frontier}\\ \hline
T&(-6,25,19)\\
AT&(-8,31,23)\\
A^2T&(-10,37,27)\\
W&(-7,26,19),\ (-7,25,20)\\
AW&(-9,32,23),\ (-9,31,24)\\
BW&(-8,31,23),\ (-8,30,24)\\
A\cU&(-7,26,21)\\
B\cU&(-6,25,21)\\
C\cU&(-4,20,18).
\end{array}
\tag{58}
\]

\begin{lemma}[Regular degree-eleven boundary]
\label{lem:regular-eleven}
At every regular graph boundary, the earliest nonzero coefficient of
\(\Jac(\cU,Q^\flat)\), for a nonconstant degree-eleven normal form
\(Q^\flat\), is not a scalar.
\end{lemma}

\begin{proof}
For \(L\ge2\), the larger \(q\)-entry dominates in each two-frontier
row of (58).  The resulting \((\beta,h+Lq)\)-labels are pairwise
distinct and distinct from every retained ordinary label.  Their
Wronskian characteristics are
\[
\begin{array}{c|c@{\qquad}c|c}
T&7L-5&AT&21L+5\\
A^2T&35L+15&W&19L+15\\
AW&33L+25&BW&16L+10\\
A\cU&14L+10&B\cU&-3L-5\\
C\cU&-22L-20&&
\end{array}
\tag{59}
\]
and none vanishes.  An exact gap audit of the retained ordinary atoms
finds seven possible vanishing characteristics.  In every case the
first nonmonomial deviation loses at most \(L\), while the next lower
label is separated by more than \(L\).

For \(L=1\), the fixed polynomial (40) leaves exactly four repeated
labels:
\[
\begin{array}{c|c}
(-4,38)&C\cU,\ AB^3C\\
(-6,46)&B\cU,\ AB^4\\
(-7,47)&A\cU,\ A^2B^3\\
(-8,54)&AT,\ BW.
\end{array}
\tag{60}
\]
For a row of \(x\)-order \(\beta\), apply
\[
\mathcal W_\beta(V)=-5H V'-\beta H'V,
\qquad H=\mathsf H(u,f(u)).
\tag{61}
\]
The differences of the normalized next-to-leading coefficients of the
two outputs in the four rows are, respectively,
\[
-\frac{481}{342},\quad
-\frac{91}{152},\quad
-\frac{2639}{1368},\quad
\frac{8968203125}{121057727177}.
\tag{62}
\]
Thus every two-dimensional map is injective, and a nonzero combination
loses at most one \(u\)-degree.  The next lower target label on each of
the four \(\beta\)-faces is at gap exactly two, so it cannot cancel the
surviving coefficient.  The two remaining singleton ordinary resonances
have first deviations of degree loss one and next-label gaps six and two.
\end{proof}

At a pole, the retained ordinary atoms and (57) have six and only six
repeated \((\beta,q)\)-labels:
\[
\begin{array}{c|c|c|c}
\text{group}&(\beta,q)&P(u)\text{ in }f^qP(u)&
B=-5q-17\beta\\ \hline
W/A^2B^3&(-7,20)&a u^{27}+b u^{25}&19\\
C\cU/A^2B^2C^2&(-4,18)&a u^{22}+b u^{20}&-22\\
B\cU/A^2B^3C&(-6,21)&a u^{27}+b u^{25}&-3\\
BW/A^2B^4&(-8,24)&a u^{32}+b u^{30}&16\\
A\cU/A^3B^2C&(-7,21)&a u^{28}+b u^{26}&14\\
AW/A^3B^3&(-9,24)&a u^{33}+b u^{31}&33.
\end{array}
\tag{63}
\]
After a nonzero monomial factor is removed, the Wronskian of a row is
\[
B\,u f'P+f\bigl(\alpha uP'-20EP\bigr).
\tag{64}
\]
If \(P(1)\ne0\), the first coefficient is nonzero because every \(B\)
in (63) is nonzero.  If \(P(1)=0\), then
\(P\) is a multiple of \(u^r(u^2-1)\), and the next coefficient is
proportional to
\[
(B\nu+\alpha)P'(1).
\tag{65}
\]
Four rows are immediately nonresonant under
\(\rho<0,\ \nu\ge2-\rho\).  The remaining equations have exactly the
two graph-compatible families
\[
\begin{aligned}
BW:\quad&
\rho=5-16k,\quad \nu=17k-5,\quad
f=c\,u^{30k-10}(u^2-1)^{17k-5},\\
A\cU:\quad&
\rho=11-14k,\quad \nu=17k-13,\quad
f=c\,u^{22k-18}(u^2-1)^{17k-13},
\end{aligned}
\qquad k\ge2,\quad c\in\C^\times.
\tag{66}
\]
Indeed, reduction of the two resonance equations modulo \(16\) and
\(14\) forces \(\rho\equiv5\pmod {16}\) and
\(\rho\equiv11\pmod {14}\), respectively; the graph inequality then
gives \(k\ge2\).  The scalar \(c\) does not change the logarithmic
derivative.

\begin{lemma}[Resonant local-jet obstruction]
\label{lem:resonant-eleven}
Neither family in (66) can cancel the first lower total-\(\gamma\)
coefficient of the exact Jacobian.
\end{lemma}

\begin{proof}
Write the top terms as
\[
x^{-5}h(u)\gamma^{17},\qquad x^\beta b(u)\gamma^n,
\qquad h(u)=\theta u^{20}.
\]
Their top-total-\(\gamma\) Jacobian is, up to a fixed nonzero monomial
factor,
\[
x^{-5+\beta}\gamma^{16+n}\mathcal L(\gamma),
\tag{67}
\]
where
\[
\begin{aligned}
\mathcal L(\gamma)={}&
hb(-5n-17\beta)\,u\gamma_u\\
&+u(17hb'-nh'b)x\gamma_x
+u(-5hb'-\beta h'b)\gamma .
\end{aligned}
\tag{68}
\]
On an \(x^I\)-sector this is a one-variable Euler equation.  Its
polynomial negative-sector solutions are exactly (66).  Linearity of
\(\mathcal L\) kills any sum of resonant intermediate sectors; a
nonresonant intermediate sector would already give an earlier isolated
coefficient.

An exact \(E=\beta+\rho q\) audit is needed here.  On the \(BW\) ray,
the complete resonant face is \(BW/A^2B^4\); on the \(A\cU\) ray it is
\(A\cU/A^3B^2C\).  No different \(q\)-group shares either face for
any \(k\ge2\).  Between that face and the first lower-total-\(\gamma\)
order, the \(BW\) ray has only the seven singleton \(q=24\) labels
\(\beta=-7,\ldots,-1\).  The \(A\cU\) ray has the three singleton
\(q=20\) labels \(\beta=-10,-9,-8\), the \(q=21\) labels
\(\beta=-6,\ldots,1\), and the duplicate \(B\cU\) atom at
\(\beta=-6\).  The singleton lemma of Section~4 closes the former
faces; the sole repeated intervening group is nonresonant because
\[
(-3)\nu+\alpha=-289k+221\ne0.
\]
Consequently a nonzero target coefficient at an intervening order would
itself be an earlier obstruction.

The first remaining sector is the \(x^0\)-sector.  With \(z=u-1\), the
graph normalization fixes its 1-jet but not its higher coefficients:
\[
f_0(u)=1-\frac{57}{34}z+z^2h_0(u).
\tag{69}
\]
At the same \(x\)-order the first lower total-\(\gamma\) seed layer
enters.  Intermediate sectors cannot enter this coefficient: replacing
a leading \(x^\rho f\) factor by one of them strictly raises its
\(x\)-order.

For the \(BW\) family, the only lower target atom at this order is
\(AT\).  The \(z^{-1},z^0\) Laurent coefficients of the normalized
residue and the leading \(AT\)-Wronskian have determinant
\[
\Delta_{BW/AT}\doteq(17k-5)(105k-32).
\tag{70}
\]
The coefficient \(h_0(1)\) cancels identically.  For the \(A\cU\)
family, the lower target space is the two-dimensional
\(W/A^2B^3\) group.  The \(z^{-1},z^0,z^1\) coefficients of the
residue and its two leading Wronskians have determinant
\[
\Delta_{A\cU/W}\doteq(17k-13)^2.
\tag{71}
\]
Both \(h_0(1)\) and \(h_0'(1)\) cancel identically.  Here \(\doteq\)
denotes equality up to a fixed nonzero rational scalar.  The two
determinants are nonzero for every \(k\ge2\), independently of every
higher coefficient in (69).

The exact polar-label audit shows that the displayed lower atoms are the
only possible interferences at these orders.  Every nonconstant monomial
of \(R_{\le2}\) starts strictly later.
\end{proof}

\begin{proof}[Proof of Theorem~\ref{thm:eleven}]
Use Lemma~\ref{lem:basis-eleven} and subtract the \(\cU\)-coefficient
modulo \(P=\cU+R_{\le2}\).  If the resulting \(Q^\flat\) is constant,
the Jacobian is zero.  Otherwise Lemma~\ref{lem:regular-eleven} handles
every regular boundary.  At a pole, (63)--(65) handle all
nonresonant groups and Lemma~\ref{lem:resonant-eleven} handles the two
families (66).  The degree-two perturbation is later in each filtration.
The isolated coefficient is either at nonzero \(x\)-order or has
nonconstant \(u\)-dependence, so the Jacobian cannot be a nonzero scalar.
\end{proof}

\section{Computational reproducibility and scope}

All exact computations use rational arithmetic.  From the repository
root, the principal checks are
\begingroup
\footnotesize
\begin{verbatim}
.venv/bin/python \
  current_context/verify_weighted_lift_degree_six_graph_projection_no_go.py
.venv/bin/python \
  current_context/verify_weighted_lift_global_minimal_boundary_face_closure.py
.venv/bin/python \
  current_context/verify_weighted_lift_nonhomogeneous_newton_vertex_closure.py
.venv/bin/python \
  current_context/verify_weighted_lift_nonhomogeneous_degree_nine_cusp_closure.py
.venv/bin/python \
  current_context/verify_weighted_lift_nonhomogeneous_degree_ten_newton_basis_closure.py
.venv/bin/python \
  current_context/verify_weighted_lift_nonhomogeneous_degree_eleven_newton_basis_closure.py
.venv/bin/python \
  current_context/verify_weighted_lift_sagbi_saturation_next_generator_audit.py
\end{verbatim}
\endgroup
The first script checks the three-dimensional Keller identity.  The
second reconstructs the seed from \(r(w)\), builds (2)--(3), checks the
exact numerator support and seed gaps, and verifies the boundary
formulas.  The third independently checks all derived Wronskian,
endpoint, lattice, cutoff, and perturbation identities.  The fourth
checks the exact degree-nine \(T\)-replacement.  The fifth checks all
four collision classes, the triangular degree-ten basis, the exact
frontiers of \(T,AT,W\), regular separation, and the exceptional polar
Wronskian.  The sixth checks all ten degree-eleven collision classes,
the augmented basis and exact frontiers, all regular and polar groups,
the exhaustive resonance parametrization, the filtration gaps, and the
two generic-zero-sector Laurent determinants.  The seventh checks the
cautionary next saturation.

The archived code version for this manuscript is
\begin{center}
\url{https://github.com/AtharvaVaidya/mathematics-research/tree/weighted-lift-newton-obstructions-v3}.
\end{center}

The theorem concerns polynomial graph embeddings into one specific
three-dimensional Keller family and a specified first target
coordinate.  It does not cover arbitrary nonlinear first-coordinate
perturbations, non-graph affine-plane embeddings, or arbitrary plane
Keller maps.  In particular it is not a proof of the plane Jacobian
conjecture.

\section*{Disclosure}

OpenAI Codex assisted with symbolic computation, proof search,
adversarial review, and drafting.  The named author is responsible for
the claims.  This manuscript should receive independent human
line-by-line verification and a broader expert novelty review before
submission.

\begin{thebibliography}{9}

\bibitem{Gallagher}
A.~Gallagher,
\emph{The Jacobian counterexample, explained},
\url{https://jacobianfun.org/jacobian-explained},
accessed 25 July 2026;
archived construction, \url{https://doi.org/10.5281/zenodo.21479195}.

\bibitem{BassConnellWright}
H.~Bass, E.~Connell, and D.~Wright,
The Jacobian conjecture: reduction of degree and formal expansion of
the inverse,
\emph{Bull.\ Amer.\ Math.\ Soc.} \textbf{7} (1982), 287--330.

\bibitem{vanDenEssen}
A.~van den Essen,
\emph{Polynomial Automorphisms and the Jacobian Conjecture},
Progress in Mathematics 190, Birkh\"auser, 2000.

\end{thebibliography}

\end{document}
